\section{Introduction}
\label{sec:introduction}

Large language models have demonstrated a striking capacity for self-directed
scientific iteration: given a task, an LLM can propose a change, execute an
experiment, observe the outcome, and decide whether to keep or discard the
change.
When repeated, this propose--execute--evaluate loop constitutes a form of
automated research~\citep{karpathy2026autoresearch}.
Instantiated for neural network hyperparameter search, we call this loop
\emph{autoresearch}.

Despite its promise, autoresearch as currently practiced has a fundamental
limitation: the search \emph{mechanism} is fixed at design time.
Every system in the literature uses a human-engineered architecture.
\citet{karpathy2026autoresearch} introduced the single-track inner loop with a
keep/discard acceptance rule.
AutoResearchClaw~\citep{aiminglab2026claw} extended it with multi-batch parallel
search.
EvoScientist~\citep{evoscientist2026} added persistent experience memory across
runs.
A human designed each improvement by reading the prior system's code,
identifying a bottleneck, and writing new code to address it.
The systems themselves cannot perform this operation.

This raises a natural question: \emph{can an outer loop perform that same
design step autonomously?}
Can the system read its own inner loop's code, identify structural
bottlenecks, generate new Python code to address them, and inject that
code at runtime---all without human involvement?

We answer this question affirmatively (\cref{fig:concept}).
We present \textbf{\methodname{}}, a bilevel framework with two
nested loops:
the \emph{inner loop} optimizes the task (proposing hyperparameter changes,
training, evaluating, keeping or discarding);
the \emph{outer loop} optimizes how the inner loop searches, by reading its
code, identifying bottlenecks, generating new Python mechanisms, and injecting
them at runtime.
Both loops use the same LLM---any improvement comes from the bilevel
\emph{architecture}, not from a more capable model.

\begin{figure}[!htbp]
  \centering
  \includegraphics[width=0.85\columnwidth]{bilevel_concept.pdf}
  \caption{Bilevel autoresearch: the inner loop optimizes the task output;
  the outer loop optimizes the inner loop's search mechanism by generating
  and injecting new Python code at runtime.}
  \label{fig:concept}
\end{figure}

We evaluate the framework on Karpathy's GPT pretraining benchmark
(\valbpb{}, 300-second training budget, RTX~5090) with a controlled
four-group ablation that isolates the contribution of each component
(details in \cref{sec:methods}).

\paragraph{Contributions.}
\begin{enumerate}
  \item We formalize autoresearch as a bilevel optimization problem in which the
        upper level optimizes the search \emph{mechanism} (runner code) and the
        lower level optimizes the task (training configuration), and demonstrate
        that the upper level can be solved by LLM code generation.
  \item We implement Level~2 mechanism research via a 4-round LLM dialogue and
        runtime code injection, establishing a concrete realization of a
        self-modifying autoresearch system.
  \item We report a controlled ablation showing that Level~2 produces a
        5$\times$ improvement over the Level-1 baseline ($-0.045 \pm 0.030$
        vs.\ $-0.009 \pm 0.002$), with Level~2 autonomously generating
        mechanisms drawn from four distinct algorithmic domains.
  \item We identify why Level~2 succeeds: the generated mechanisms (Tabu
        Search, Multi-Scale Bandit, Orthogonal Exploration) forced exploration
        of \texttt{TOTAL\_BATCH\_SIZE} \emph{reduction}, a direction that the
        LLM's default search path systematically avoids due to an implicit bias
        toward larger batches.
\end{enumerate}

\Cref{sec:related} reviews related work.
\Cref{sec:methods} describes the framework in detail.
\Cref{sec:results} presents ablation results.
\Cref{sec:discussion} interprets hypotheses, failure modes, and limitations.
\Cref{sec:conclusion} concludes.
